\chapter{Mechanistic Interpretability: Estrazione e Manipolazione Latente}
\chaptermark{Mechanistic Interpretability}

Questo capitolo esplora l'analisi e la manipolazione dei modelli linguistici attraverso la lente della Mechanistic Interpretability. L'obiettivo primario è superare il paradigma puramente osservativo tipico degli attacchi Black-Box, limitato allo studio della correlazione tra input e output, per penetrare all'interno della rete e indagarne i meccanismi causali. Attraverso operazioni mirate, si intende dimostrare come l'architettura profonda del Transformer giochi un ruolo cruciale non solo nella suscettibilità alle minacce esterne, ma soprattutto nella genesi stessa del processo decisionale.\\

In questo contesto, l'indagine si concentra sull'attacco White-Box noto come Activation Steering. Grazie al suo intervento diretto sulla geometria dello spazio latente e alla sua applicazione chirurgica su specifici layer della rete, questa tecnica trascende il semplice scopo offensivo, configurandosi come un potente strumento diagnostico per mappare le dinamiche interne del modello.\\

L'analisi condotta in queste pagine assume pertanto una triplice valenza. Dal punto di vista offensivo, permette di studiare come raffinare l'attacco, rendendolo più silenzioso ed efficiente. Dal punto di vista difensivo, apre la strada alla valutazione della resilienza della rete e alla neutralizzazione delle perturbazioni tramite interventi di compensazione, noti in letteratura come Positive Steering \cite{Wendlinger2026}. Infine, da un punto di vista strettamente architetturale, lo studio degli stati interni consente di dissezionare il flusso computazionale del modello, isolando il contributo delle singole partizioni di layer al fine di identificare l'esatto punto di rottura, ovvero l'anello più debole nella catena logica della rete.
Infine l'activation steering permette di avere una base di riferimento nello spazio latente: poter mappare il vettore sicuro-vulnerabile attraverso l'osservazione degli hidden states getta le basi per un'analisi e paragone con gli attacchi white-box

\section{Estrazione dei vettori}
Nella prospettiva della \textit{Mechanistic Interpretability}, per poter calcolare la direzione vettoriale della vulnerabilità non è sufficiente osservare l'output finale, ma è necessario mappare l'intero processo cognitivo interno della rete. È stato quindi implementata un'infrastruttura di monitoraggio capace di estrarre e registrare gli stati latenti che fluiscono attraverso il \textit{Residual Stream} del Transformer durante l'elaborazione degli snippet di codice. Questa operazione richiede di \enquote{aprire} la \textit{black-box} del modello e catturarne lo stato in tempo reale, senza tuttavia alterarne l'architettura originaria o i pesi.

All'interno del Residual Stream di un Transformer, l'informazione fluisce sotto forma di tensori, ovvero matrici multidimensionali di dati. Ad ogni layer la rete proietta ciascun token di testo in uno spazio iperdimensionale, codificandolo attraverso un vettore nello spazio latente. I singoli valori numerici che compongono questo vettore prendono il nome di attivazioni.

L'estrazione massiva dei dati a scopo diagnostico è stata condotta sfruttando le funzionalità native dell'ecosistema HuggingFace. Durante il forward pass, è stato abilitato il parametro output\_hidden\_states=True. Questa scelta architetturale istruisce il modello a preservare e restituire l'intera traiettoria computazionale in una singola inferenza, fornendo accesso simultaneo ai tensori di attivazione di tutti i layer del Transformer, omettendo unicamente il layer di embedding iniziale in quanto considerato di natura lessicale e inutile ai fini dello studio.

\begin{minted}[linenos, frame=lines, framesep=2mm, fontsize=\footnotesize]{python}
with torch.no_grad():
    outputs = model(**inputs, output_hidden_states=True)

hidden_states = outputs.hidden_states[1:] 

for layer_idx in range(num_layers):
    vettore_ultimo_token = hidden_states[layer_idx][0, -1, :].detach().cpu().numpy()

    if target == 'Vulnerabile':
        attivazioni_vulnerabili[layer_idx].append(vettore_ultimo_token)
    elif target == 'Sicuro':
        attivazioni_sicuri[layer_idx].append(vettore_ultimo_token)

\end{minted}

Un aspetto metodologico cruciale di questa fase riguarda la selezione delle coordinate tensoriali da estrarre. L'output di un layer Transformer è un tensore tridimensionale con dimensioni [Batch, Sequenza, Neuroni]. Per gli scopi di questo studio, l'estrazione si è concentrata esclusivamente sulle attivazioni dell'ultimo token della sequenza, indicizzato come [0, -1, :]. Nei modelli autoregressivi decoder-only, il meccanismo di causal masking impedisce ai token di guardare al futuro, di conseguenza l'ultimo token elaborato in fase di prefill è l'unico ad avere applicato l'attenzione sull'intero prompt di ingresso, aggregando in sé l'intera rappresentazione semantica e logica del codice appena letto.

Una volta ottenuti, i vettori estratti vengono associati ai rispettivi metadati, id dello snippet, ground truth di sicurezza. Questa raccolta di attivazioni latenti costituisce il materiale grezzo su cui verranno successivamente calcolati i centroidi per la generazione dello Steering Vector per quel layer.

\begin{minted}[linenos, frame=lines, framesep=2mm, fontsize=\footnotesize]{python}
 media_vuln = np.mean(np.stack(attivazioni_vulnerabili[layer_idx]), axis=0)
            media_sicuro = np.mean(np.stack(attivazioni_sicuri[layer_idx]), axis=0)
            vettore_steering = media_vuln - media_sicuro
\end{minted}

I centroidi vengono calcolati esattamente come illustrato nell'attacco di steering: viene eseguita la media dei vettori corrispondenti ai codici vulnerabili e sicuri ottenendo così due vettori. 
Tramite la differenza vettoriale tra i centroidi si ottiene lo Steering Vector del modello: reiterando questo procedimento ad ogni layer si estrae l'intero asse di riferimento, ricavando per ciascun livello un vettore orientato verso il centroide vulnerabile.



\section{Activation Steering}
L'Activation Steering rappresenta una tecnica avanzata di manipolazione white-box che interviene direttamente nello spazio latente del modello. Alterando le attivazioni interne durante la fase di inferenza, è possibile deviare la traiettoria generativa dell'LLM verso concetti o comportamenti desiderati. L'efficacia di questa operazione dipende fortemente da parametri finemente modulabili, in particolare il layer di iniezione, il punto in cui avviene la perturbazione, e il moltiplicatore, l'intensità della perturbazione stessa. 

Eseguire uno studio di ablazione su questi parametri permette di raggiungere un duplice obiettivo: da un lato, ottimizzare l'attacco per studiare i limiti e la resilienza del modello, individuando le vulnerabilità intrinseche e i layer più deboli della sua architettura Transformer; dall'altro, mappare l'andamento logico del modello nello spazio latente, osservando come la semantica e la logica si degradano o si trasformano durante l'elaborazione nella rete.

\subsection{Ricerca del Moltiplicatore}
Il moltiplicatore agisce come un fattore di scala per il vettore di steering innescato. Un valore troppo basso risulterebbe inefficace, venendo assorbito dalle normali computazioni del modello, mentre un valore troppo alto rischierebbe di distruggere la coerenza del testo generato portando a gibberish.

La ricerca del moltiplicatore ideale per ciascun modello prevede l'esecuzione di uno studio di ablazione fissato al layer medio, considerato in letteratura il punto di transizione semantica più malleabile. In questa fase preliminare, l'attacco di Activation Steering è stato testato con la seguente progressione di moltiplicatori: 3, 5, 8, 10, 15.

L'indagine si configura come la ricerca di un trade-off tra il rateo di successo dell'attacco (ASR) e la percentuale di errori o di testo senza senso. Come avviene per altre tecniche come quella del Logit Bias, è teoricamente possibile forzare il modello per ottenere un ASR quasi perfetto; tuttavia, se questo successo è accompagnato da un'elevata confusione e dalla perdita di logica nella semantica, l'output diventa inutilizzabile. Un attacco white-box efficace deve eludere i filtri di sicurezza mantenendo intatte le capacità di instruction-following e di articolazione del linguaggio del modello. I risultati di questa prima fase sono riportati nella Tabella \ref{tab:valutazione_moltiplicatore}.


\begin{table}[H]
    \centering
    \resizebox{\textwidth}{!}{
    \begin{tabular}{l c c c c}
        \toprule
        \textbf{Modelli} & \textbf{Layer} & \textbf{Moltiplicatore} & \textbf{ASR\_percentuale} & \textbf{Errori\_percentuale} \\
        \midrule
        Qwen/Qwen2.5-7B-Instruct                 & 14  & 3 & 7.45  & 0 \\
        Qwen/Qwen2.5-7B-Instruct                 & 14  & 5 & 13.83  & 0\\
        Qwen/Qwen2.5-7B-Instruct                 & 14  & 8 & 28.72  & 0 \\
        Qwen/Qwen2.5-7B-Instruct                 & 14  & 10 & 35.11  & 0 \\
        Qwen/Qwen2.5-7B-Instruct                 & 14  & 15 & 53.19  & 0 \\
        \midrule
        Qwen/Qwen2.5-Coder-7B-Instruct           & 14  & 3 & 6.17  & 0 \\
        Qwen/Qwen2.5-Coder-7B-Instruct           & 14  & 5 & 9.88  & 0 \\
        Qwen/Qwen2.5-Coder-7B-Instruct           & 14  & 8 & 13.58 & 0 \\
        Qwen/Qwen2.5-Coder-7B-Instruct           & 14  & 10 & 14.81  & 0 \\
        Qwen/Qwen2.5-Coder-7B-Instruct           & 14  & 15 & 24.69  & 0 \\
        \midrule
        meta-llama/Llama-3.1-8B-Instruct         & 16  & 3 & 69.04 & 1.24 \\
        meta-llama/Llama-3.1-8B-Instruct         & 16  & 5 & 91.02  & 8.67 \\
        meta-llama/Llama-3.1-8B-Instruct         & 16  & 8 & 47.99  & 52.01 \\
        meta-llama/Llama-3.1-8B-Instruct         & 16  & 10 & 3.41  & 96.59 \\
        meta-llama/Llama-3.1-8B-Instruct         & 16  & 15 & 0.00  & 100.00 \\
        \midrule
        codellama/CodeLlama-7b-Instruct-hf       & 16 & 3 & 22.11  & 0.50 \\
        codellama/CodeLlama-7b-Instruct-hf       & 16 & 5 & 31.66  & 0.50 \\
        codellama/CodeLlama-7b-Instruct-hf       & 16 & 8 & 53.27 & 0.50 \\
        codellama/CodeLlama-7b-Instruct-hf       & 16 & 10 & 62.31  & 1.01 \\
        codellama/CodeLlama-7b-Instruct-hf       & 16 & 15 & 76.88  & 1.01 \\
        \midrule
        deepseek-ai/DeepSeek-R1-Distill-Qwen-7B  & 14  & 3 & 11.41  & 7.22 \\
        deepseek-ai/DeepSeek-R1-Distill-Qwen-7B  & 14  & 5 & 15.97  & 4.56 \\
        deepseek-ai/DeepSeek-R1-Distill-Qwen-7B  & 14  & 8 & 15.59  & 6.84 \\
        deepseek-ai/DeepSeek-R1-Distill-Qwen-7B  & 14  & 10 & 17.87  & 5.32 \\
        deepseek-ai/DeepSeek-R1-Distill-Qwen-7B  & 14  & 15 & 15.21  & 5.70 \\
        \midrule
        deepseek-ai/deepseek-coder-6.7b-instruct & 16  & 3 & 28.57  & 0 \\
        deepseek-ai/deepseek-coder-6.7b-instruct & 16  & 5 & 57.14  & 0 \\
        deepseek-ai/deepseek-coder-6.7b-instruct & 16  & 8 & 85.71  & 0 \\
        deepseek-ai/deepseek-coder-6.7b-instruct & 16  & 10 & 85.71  & 0 \\
        deepseek-ai/deepseek-coder-6.7b-instruct & 16  & 15 & 85.71  & 0 \\
        \bottomrule
    \end{tabular}
    }
    \vspace{0.2cm}
    \caption{Ricerca del moltiplicatore}
    \label{tab:valutazione_moltiplicatore}
\end{table}

Analizzando i dati, emerge un comportamento distintivo della famiglia di modelli Qwen: a differenza di Llama, che subisce un rapido collasso semantico all'aumentare della perturbazione, i modelli Qwen registrano un errore percentuale costantemente pari a zero accompagnato da un ASR in lenta ma costante crescita. Questa spiccata robustezza topologica ha reso necessaria l'esecuzione di un secondo studio di ablazione mirato esclusivamente ai modelli di Alibaba, applicando moltiplicatori decisamente più elevati (20, 25, 30 e 40) per forzare il superamento dei filtri.

\begin{table}[H]
    \centering
    \resizebox{\textwidth}{!}{
    \begin{tabular}{l c c c c}
        \toprule
        \textbf{Modelli} & \textbf{Layer} & \textbf{Moltiplicatore} & \textbf{ASR\_percentuale} & \textbf{Errori\_percentuale} \\
        \midrule
        Qwen/Qwen2.5-7B-Instruct                 & 14  & 20 & 73.40  & 0 \\
        Qwen/Qwen2.5-7B-Instruct                 & 14  & 25 & 89.36  & 0\\
        Qwen/Qwen2.5-7B-Instruct                 & 14  & 30 & 94.68  & 0 \\
        Qwen/Qwen2.5-7B-Instruct                 & 14  & 40 & 98.94  & 0 \\
        \midrule
        Qwen/Qwen2.5-Coder-7B-Instruct           & 14  & 20 & 29.63  & 0 \\
        Qwen/Qwen2.5-Coder-7B-Instruct           & 14  & 25 & 30.86  & 0 \\
        Qwen/Qwen2.5-Coder-7B-Instruct           & 14  & 30 & 35.80 & 0 \\
        Qwen/Qwen2.5-Coder-7B-Instruct           & 14  & 40 & 44.44  & 0 \\
      
        \bottomrule
    \end{tabular}
    }
    \vspace{0.2cm}
    \caption{Ricerca del moltiplicatore aggiunta}
    \label{tab:valutazione_moltiplicatore_aggiunta}
\end{table}

Come si evince dalla Tabella \ref{tab:valutazione_moltiplicatore_aggiunta}, l'ASR dei modelli Qwen continua a crescere mantenendo nullo il tasso di errore. Sebbene valori estremi, come 40, portino al rateo di successo più alto in assoluto, deformare in modo così aggressivo il vettore nello spazio latente rischia di introdurre instabilità non immediatamente percepibili dalla metrica degli errori, compromettendo l'affidabilità generale del modello, E' preferibile quindi adottare una metodologia più conservativa.

Per garantire un'applicazione adeguata, coerente e bilanciata dell'operazione di steering, sono stati selezionati definitivamente i seguenti moltiplicatori per le fasi successive della ricerca:

\begin{itemize}[noitemsep]
    \item Qwen/Qwen2.5-7B-Instruct: 20
    \item Qwen/Qwen2.5-Coder-7B-Instruct: 30
    \item meta-llama/Llama-3.1-8B-Instruct: 3
    \item codellama/CodeLlama-7b-Instruct-hf: 15
    \item deepseek-ai/DeepSeek-R1-Distill-Qwen-7B: 10
    \item deepseek-ai/deepseek-coder-6.7b-instruct: 8
\end{itemize}

\subsection{Layer Sweep}
Esattamente come il moltiplicatore, la scelta del layer di iniezione ha un ruolo fondamentale per massimizzare l'operazione di attacco: se eseguito troppo precocemente in un layer la rete è ancora impegnata nell'analisi lessicale del testo e tende ad assorbire o correggere l'anomalia vettoriale.
Al contrario l'iniezione effettuata tardivamente il modello ha già consolidato il proprio verdetto logico e la perturbazione risulta ininfluente.

E' stato eseguito un layer sweep su tutti i layer di tutti i modelli con i parametri precedentemente ottenuti. Attraverso un ciclo viene eseguita un'operazione massiva: applicando lo steering ad un layer e cambiandolo ad ogni iterazione con quello successivo è possibile tracciare l'andamento dell'efficacia dell'iniezione.

\begin{figure}[H]
    \centering
    
    \begin{subfigure}[b]{0.45\textwidth}
        \centering
        \includegraphics[width=\textwidth]{immagini/Layer Sweep/Qwen_Qwen2.5-7B-Instruct.png} 
        \caption{Risultati per Qwen2.5-7B}
        \label{fig:qwen_7b}
    \end{subfigure}
    \hfill 
    \begin{subfigure}[b]{0.45\textwidth}
        \centering
        \includegraphics[width=\textwidth]{immagini/Layer Sweep/Qwen_Qwen2.5-Coder-7B-Instruct.png}
        \caption{Risultati per Qwen2.5-Coder}
        \label{fig:qwen_coder}
    \end{subfigure}
    
    \vspace{0.5cm} 
    
    \begin{subfigure}[b]{0.45\textwidth}
        \centering
        \includegraphics[width=\textwidth]{immagini/Layer Sweep/meta-llama_Llama-3.1-8B-Instruct.png}
        \caption{Risultati per Llama-3.1-8B}
        \label{fig:llama_3_1}
    \end{subfigure}
    \hfill
    \begin{subfigure}[b]{0.45\textwidth}
        \centering
        \includegraphics[width=\textwidth]{immagini/Layer Sweep/codellama_CodeLlama-7b-Instruct-hf.png}
        \caption{Risultati per CodeLlama-7b}
        \label{fig:codellama}
    \end{subfigure}
    
    \vspace{0.5cm} 
    
    \begin{subfigure}[b]{0.45\textwidth}
        \centering
        \includegraphics[width=\textwidth]{immagini/Layer Sweep/deepseek-ai_DeepSeek-R1-Distill-Qwen-7B.png}
        \caption{Risultati per DeepSeek-R1}
        \label{fig:deepseek_r1}
    \end{subfigure}
    \hfill
    \begin{subfigure}[b]{0.45\textwidth}
        \centering
        \includegraphics[width=\textwidth]{immagini/Layer Sweep/deepseek-ai_deepseek-coder-6.7b-instruct.png}
        \caption{Risultati per Deepseek-Coder}
        \label{fig:deepseek_coder}
    \end{subfigure}
    
    \vspace{0.3cm}
    \caption{Layer Sweep: andamento dell'attacco per i sei modelli analizzati.}
    \label{fig:griglia_layer_sweep}
\end{figure}

Come è possibile osservare dai grafici, l'attacco ottiene un alto ASR generalmente nel layer intermedi dei modelli, confermando la teoria del partizionamento dei layer dei modelli: il basso successo della manipolazione nei layer iniziali e in quelli finali confermano l'effettiva importanza nell'elaborazione della risposta del modello. Viene posto un forte accento sui layer intermedi, responsabili del consolidamento semantico e della formazione dei concetti logici complessi come il ragionamento sulla vulnerabilità del codice, prima che questi vengano tradotti in specifici token di output.
Tramite questa operazione è possibile ottenere il layer più influenzabile: il livello che presenta il rateo di successo dello steering più alto, indica la posizione più sensibile a livello architetturale del modello.\ref{tab:esito_layer_sweep}

\begin{table}[H]
    \centering
    \resizebox{\textwidth}{!}{
    \begin{tabular}{l c c c c}
        \toprule
        \textbf{Modelli} & \textbf{Layer} & \textbf{Moltiplicatore} & \textbf{ASR\_percentuale} & \textbf{Errori\_percentuale} \\
        \midrule
        Qwen/Qwen2.5-7B-Instruct                 & 18  & 20 & 100  & 0 \\
        Qwen/Qwen2.5-Coder-7B-Instruct                & 18  & 30 & 100  & 0\\
        meta-llama/Llama-3.1-8B-Instruct               & 15  & 3 & 89.16  & 1.86 \\
        codellama/CodeLlama-7b-Instruct-hf               & 13  & 15 & 97.49  & 0.50 \\
        deepseek-ai/DeepSeek-R1-Distill-Qwen-7B           & 19  & 10 & 19.01  & 3.04 \\
        deepseek-ai/deepseek-coder-6.7b-instruct         & 16  & 8 & 85.71  & 0 \\
        \bottomrule
    \end{tabular}
    }
    \vspace{0.2cm}
    \caption{Esito ricerca del layer più sensibile}
    \label{tab:esito_layer_sweep}
\end{table}

Attraverso questi risultati è possibile ottenere un "\textit{Golden Layer}" identificativo per la parte di maggior computazione del modello: calcolando la percentuale di profondità del layer sensibile per ogni modello e poi facendo la media di questi si ottiene l'altezza media dei layer per cui, in caso non si potesse eseguire un layer sweep o l'inaccessibilità al modello, è possibile applicare attacchi, o operazioni, nella parte più sensibile.
Nello specifico con questi sei modelli si può notare come solo per il modello generale di llama il layer più sensibile sia quello centrale, mentre per gli altri modelli è leggermente più spostato ma sempre nella porzione intermedia dell'architettura dei modelli.
Considerando quindi i dati ottenuti è possibile approssimare il golden layer a circa il 55\% dell'altezza del modello.

Di seguito il calcolo della profondità relativa per ciascun modello:
\begin{itemize}[noitemsep]
    \item \textbf{Qwen/Qwen2.5-7B-Instruct:} $18/28 = 0{,}643 \rightarrow 64{,}3\%$
    \item \textbf{Qwen/Qwen2.5-Coder-7B-Instruct:} $18/28 = 0{,}643 \rightarrow 64{,}3\%$
    \item \textbf{meta-llama/Llama-3.1-8B-Instruct:} $15/32 = 0{,}469 \rightarrow 46{,}9\%$
    \item \textbf{codellama/CodeLlama-7b-Instruct-hf:} $13/32 = 0{,}406 \rightarrow 40{,}6\%$
    \item \textbf{deepseek-ai/DeepSeek-R1-Distill-Qwen-7B:} $19/28 = 0{,}679 \rightarrow 67{,}9\%$
    \item \textbf{deepseek-ai/deepseek-coder-6.7b-instruct:} $16/32 = 0{,}500 \rightarrow 50{,}0\%$
\end{itemize}

Facendo la media delle percentuali calcolate si ottiene:
$$ \frac{64{,}3+64{,}3+46{,}9+40{,}6+67{,}9+50{,}0}{6} = \frac{334}{6} = 55{,}66\% $$

Ai fini dello studio verranno adottati i layer più sensibili senza approssimare alla percentuale di 56\%, in uno studio in cui la ricerca del layer non fosse possibile sarebbe necessario adottare il livello di profondità ottenuto dalla media precedente.

